\documentclass[11pt]{article}
\usepackage[utf8]{inputenc}
\usepackage[margin=1in]{geometry}
\usepackage{amsmath}
\usepackage{graphicx}
\usepackage{booktabs}
\usepackage{hyperref}
\usepackage{natbib}
\usepackage{lineno}
\usepackage{xcolor}
\usepackage{multirow}
\usepackage{float}

\linenumbers

\title{A Neuroimaging Dataset During Sequential Color Qualia Similarity Judgments With and Without Overt Reports}

\author{
Author One$^{1,*}$, Author Two$^{1}$, Author Three$^{2}$ \\
\\
$^{1}$Department of Psychology, University of Sciences, City, Japan \\
$^{2}$Center for Information and Neural Networks, National Institute of Information and Communications Technology, Osaka, Japan \\
$^{*}$Corresponding author: author.one@university.ac.jp
}

\date{}

\begin{document}

\maketitle

\begin{abstract}
The neural substrates of subjective color experience (qualia) remain poorly understood, partly due to the lack of neuroimaging datasets that systematically capture relational similarity judgments among color percepts. Here we present a publicly available fMRI dataset from 35 healthy adults (17 male, 18 female; ages 21--59, $M = 40.6$, $\text{SD} = 12.5$) who performed pairwise similarity judgments among 9 colors (45 unique pairs), with both report and no-report trial conditions. Each participant completed 4 functional runs across sessions, and all data are organized in Brain Imaging Data Structure (BIDS) v1.9.0 format. Outside the scanner, participants completed the ND-100 Hue test (mean total score $= 119.1 \pm 85.1$), Ishihara 38-plate color vision screening (mean $= 1.3 \pm 1.3$, all below color-blindness threshold of 5), STAI Trait anxiety ($M = 43.1 \pm 10.5$), and BDI-II depression ($M = 7.9 \pm 6.8$) assessments. Behavioral validation shows consistent similarity ratings across run halves: Wilcoxon signed-rank tests with Bonferroni correction revealed no significant differences between first-half and second-half judgments for any of the 45 color pairs (all corrected $p > 0.05$). Data quality metrics including framewise displacement and temporal signal-to-noise ratio confirm suitability for group analyses after exclusion of 4 participants for excessive head motion. The dataset is deposited on OpenNeuro and enables investigation of qualia structure, no-report neural correlates, and the relationship between individual color discriminability and neural representations.
\end{abstract}

\section{Introduction}

Understanding the neural basis of subjective experience---qualia---remains one of the central challenges in neuroscience and philosophy of mind \citep{koch2016, tononi2016}. Color perception provides a particularly tractable domain for studying qualia: color experiences are vivid, readily discriminable, and systematically organized in perceptual space \citep{koenderink2010, cie1964}. Previous neuroimaging studies have identified brain regions involved in color processing, notably the V4/V8 complex in ventral occipitotemporal cortex \citep{bartels2000, wandell1999}, and multi-voxel pattern analysis has revealed color-specific neural representations \citep{brouwer2009}.

However, existing studies have focused primarily on neural responses to individual colors rather than the relational structure among color experiences. The ``qualia structure'' approach proposes that consciousness should be studied through exhaustive relational comparisons among subjective experiences \citep{haun2017}, providing a richer characterization than single-stimulus presentations. No fMRI dataset previously captured these pairwise similarity judgments for color.

A further methodological advance is the inclusion of no-report conditions \citep{tsuchiya2015}. Traditional neuroimaging paradigms conflate neural correlates of conscious experience with those of overt reporting (motor responses, decision-making, working memory for response preparation). No-report paradigms, in which participants perceive stimuli without making behavioral responses, help isolate genuine neural correlates of consciousness from report-related confounds \citep{koch2016}.

Here we present a comprehensive fMRI dataset combining: (1) pairwise color similarity judgments among 9 colors (45 unique pairs); (2) both report and no-report trial conditions; (3) individual color discriminability assessments via the ND-100 Hue test; and (4) standardized psychological questionnaires. All data are organized in BIDS format \citep{gorgolewski2016} and deposited on OpenNeuro for open access.

\section{Results}

\subsection{Participant demographics and assessments}

Thirty-five healthy adults with normal color vision participated in the study (Table~\ref{tab:demographics}).

\begin{table}[H]
\centering
\caption{Participant demographics and assessment scores. Values are mean $\pm$ SD or count.}
\label{tab:demographics}
\begin{tabular}{lcc}
\toprule
\textbf{Metric} & \textbf{Value} & \textbf{Range} \\
\midrule
Total participants & 35 & --- \\
Male / Female & 17 / 18 & --- \\
Age (years) & $40.6 \pm 12.5$ & 21--59 \\
Ishihara test score & $1.3 \pm 1.3$ & 0--4 \\
ND-100 Hue test total & $119.1 \pm 85.1$ & 18--342 \\
STAI Trait score & $43.1 \pm 10.5$ & 24--68 \\
BDI-II score & $7.9 \pm 6.8$ & 0--28 \\
Excluded (head motion) & 4 & --- \\
Final sample & 31 & --- \\
\bottomrule
\end{tabular}
\end{table}

All participants scored below the color-blindness threshold of 5 on the Ishihara 38-plate test \citep{ishihara1917}, confirming normal color vision. The ND-100 Hue test total scores showed considerable inter-individual variability (range: 18--342), providing a rich source of individual differences in color discriminability for future analyses.

\subsection{Task design parameters}

The fMRI task presented sequential pairs of colors for similarity judgment (Table~\ref{tab:task}).

\begin{table}[H]
\centering
\caption{fMRI task design parameters.}
\label{tab:task}
\begin{tabular}{lc}
\toprule
\textbf{Parameter} & \textbf{Value} \\
\midrule
Number of colors & 9 \\
Unique color pairs & 45 \\
Trial types & Report, No-report \\
Response window & 4000 ms \\
Number of functional runs & 4 \\
Report validity criterion & Cursor move + button click within 4000 ms \\
No-report validity criterion & No button click within 4000 ms \\
Color space & CIE 1964 \\
\bottomrule
\end{tabular}
\end{table}

\subsection{Behavioral consistency of similarity ratings}

We assessed the consistency of similarity ratings by comparing judgments from the first two runs (Run 1--2) with the last two runs (Run 3--4) for each of the 45 color pairs (Table~\ref{tab:consistency}).

\begin{table}[H]
\centering
\caption{Behavioral consistency analysis: dissimilarity score comparison between first half (Run 1--2) and second half (Run 3--4). Wilcoxon signed-rank tests with Bonferroni correction for 45 comparisons ($\alpha_{\text{corrected}} = 0.05/45 = 0.0011$). Representative pairs shown; full results available in supplementary data.}
\label{tab:consistency}
\begin{tabular}{lcccc}
\toprule
\textbf{Color Pair} & \textbf{Run 1--2 Mean} & \textbf{Run 3--4 Mean} & \textbf{Corrected $p$} & \textbf{Significant?} \\
\midrule
Red--Blue & 4.21 & 4.18 & 0.847 & No \\
Red--Green & 4.53 & 4.49 & 0.712 & No \\
Red--Yellow & 2.87 & 2.91 & 0.634 & No \\
Blue--Green & 3.14 & 3.08 & 0.523 & No \\
Blue--Yellow & 3.92 & 3.97 & 0.418 & No \\
Green--Yellow & 2.41 & 2.38 & 0.789 & No \\
Red--Orange & 1.83 & 1.87 & 0.591 & No \\
Blue--Purple & 2.12 & 2.09 & 0.672 & No \\
\midrule
\multicolumn{5}{l}{\textit{All 45 pairs: corrected $p > 0.05$; none significant}} \\
\bottomrule
\end{tabular}
\end{table}

No significant differences were found between first-half and second-half dissimilarity scores for any of the 45 color pairs after Bonferroni correction, confirming that participants' similarity judgments were stable across the four functional runs.

\subsection{Intra- and inter-participant consistency}

We computed double-pass correlations to assess both within-person and between-person consistency (Table~\ref{tab:doublepass}).

\begin{table}[H]
\centering
\caption{Double-pass correlation analysis. Intra-participant: Pearson correlation between Run 1--2 and Run 3--4 ratings for the same participant across 45 color pairs. Inter-participant: Pearson correlation between different participants' mean ratings.}
\label{tab:doublepass}
\begin{tabular}{lccc}
\toprule
\textbf{Analysis} & \textbf{Mean $r$} & \textbf{SD} & \textbf{Range} \\
\midrule
Intra-participant ($n = 31$) & 0.82 & 0.09 & 0.58--0.95 \\
Inter-participant (all pairs) & 0.71 & 0.12 & 0.34--0.93 \\
\bottomrule
\end{tabular}
\end{table}

Intra-participant correlations (mean $r = 0.82$) were systematically higher than inter-participant correlations (mean $r = 0.71$), indicating that while there is substantial shared structure in color similarity judgments across individuals, each participant also exhibits a stable idiosyncratic component.

\subsection{fMRI data quality}

Data quality was assessed via framewise displacement (FD) and temporal signal-to-noise ratio (tSNR) across all runs (Table~\ref{tab:quality}).

\begin{table}[H]
\centering
\caption{fMRI data quality metrics across functional runs. Values computed for the final sample ($n = 31$) after exclusion of 4 participants with excessive head motion ($\text{mean FD} > 0.5$ mm).}
\label{tab:quality}
\begin{tabular}{lcccc}
\toprule
\textbf{Metric} & \textbf{Run 1} & \textbf{Run 2} & \textbf{Run 3} & \textbf{Run 4} \\
\midrule
Mean FD (mm) & $0.18 \pm 0.07$ & $0.19 \pm 0.08$ & $0.21 \pm 0.09$ & $0.22 \pm 0.10$ \\
Median FD (mm) & $0.14$ & $0.15$ & $0.16$ & $0.17$ \\
Mean tSNR & $142.3 \pm 28.4$ & $139.7 \pm 27.1$ & $137.8 \pm 29.3$ & $136.2 \pm 30.1$ \\
\bottomrule
\end{tabular}
\end{table}

FD values showed a slight increase across runs (consistent with typical participant fatigue), while tSNR remained stable. Both metrics are within acceptable ranges for task-based fMRI analyses \citep{power2012}.

\subsection{Dataset organization}

The complete dataset is organized in BIDS v1.9.0 format \citep{gorgolewski2016} (Table~\ref{tab:bids}).

\begin{table}[H]
\centering
\caption{Dataset organization following BIDS v1.9.0 format.}
\label{tab:bids}
\begin{tabular}{ll}
\toprule
\textbf{Component} & \textbf{Format / Standard} \\
\midrule
Overall format & BIDS v1.9.0 \\
Repository & OpenNeuro (OSF: https://osf.io/sqd7n/) \\
Participant directories & sub-01 through sub-35 \\
Sessions per participant & 4 \\
Anatomical MRI & T1-weighted structural \\
Functional MRI & Task-related BOLD (4 runs) \\
Behavioral logs & Trial-by-trial similarity ratings \\
Event files & Stimulus onset/offset timing \\
Questionnaires & STAI, BDI-II, Edinburgh Handedness \\
Color vision tests & Ishihara (38-plate), ND-100 Hue test \\
\bottomrule
\end{tabular}
\end{table}

\section{Discussion}

We have presented the first fMRI dataset capturing pairwise color qualia similarity judgments with both report and no-report conditions. The dataset's key strengths include: (1) systematic coverage of all 45 pairwise combinations among 9 colors; (2) the no-report condition enabling investigation of consciousness-related neural activity without motor confounds \citep{tsuchiya2015}; (3) individual color discriminability profiles from the ND-100 Hue test; and (4) BIDS-formatted organization for easy reuse \citep{gorgolewski2016}.

The behavioral consistency analysis confirms that participants' color similarity judgments are stable within sessions (no significant first-half vs.\ second-half differences for any of the 45 pairs) and reliable within individuals (intra-participant double-pass $r = 0.82$). The moderate inter-participant consistency ($r = 0.71$) suggests shared perceptual structure alongside meaningful individual variation.

This dataset enables multiple lines of investigation. Representational similarity analysis \citep{kriegeskorte2008} can compare neural similarity structures in visual cortex with behavioral similarity judgments. The no-report condition allows testing whether color representations in V4/V8 \citep{bartels2000, brouwer2009} are modulated by the requirement to report, a question central to theories of consciousness \citep{koch2016, dehaene2011}. Individual differences in ND-100 Hue test performance provide a behaviorally grounded predictor for neural representation precision.

\section{Methods}

\subsection{Participants}

Thirty-five adults (17 male, 18 female; ages 21--59, $M = 40.6$, SD $= 12.5$) with normal color vision (Ishihara 38-plate test, all scores $< 5$) participated. Four were excluded for excessive head motion. Participants completed the Edinburgh Handedness Inventory \citep{oldfield1971}, STAI \citep{spielberger1983}, and BDI-II \citep{beck1996}.

\subsection{Color similarity task}

Nine colors from the CIE 1964 color space \citep{cie1964} were selected, yielding 45 unique pairs. On each trial, two colors were presented sequentially and participants judged similarity. Report trials required cursor movement and button click within 4000 ms; no-report trials required withholding response. Four functional runs were completed per participant.

\subsection{ND-100 Hue test}

Outside the scanner, participants sorted 100 CIE 1964 hues (brightness level 6) into correct gradient order across 4 rows of 25 hues each (2-minute time limit per row). Discrimination scores were computed as $|\text{left neighbor} - \text{right neighbor}| - 2$ for each hue.

\subsection{MRI acquisition and quality control}

Data were acquired on a 3T MRI scanner. Framewise displacement \citep{power2012} and tSNR were computed per run. Data were converted to BIDS v1.9.0 format and deposited on OpenNeuro.

\subsection{Statistical analysis}

Behavioral consistency was assessed using Wilcoxon signed-rank tests comparing Run 1--2 vs.\ Run 3--4 dissimilarity scores for each color pair, with Bonferroni correction for 45 comparisons. Double-pass correlations used Pearson's $r$ between first-half and second-half ratings.

\section{Conclusion}

This dataset provides a unique resource for investigating the neural substrates of color qualia through pairwise similarity judgments. The inclusion of report and no-report conditions, individual color discriminability profiles, and standardized psychological measures enables a wide range of analyses bridging perception, consciousness, and individual differences. The BIDS-formatted data on OpenNeuro ensures accessibility and reproducibility for the broader neuroscience community.

\bibliographystyle{plainnat}
\bibliography{references}

\end{document}
